% paper/sections/04f_uccd6.tex
%
% §4.5  UCCD6：6 次節点中心 CCD ＋ 次数保存ハイパー粘性
%
% ═══════════════════════════════════════════════════════════════════════════════
\subsection{\texorpdfstring{{\UCCD[6]}：次数保存ハイパー粘性付き節点中心 CCD}{UCCD6：次数保存ハイパー粘性付き節点中心 CCD}}
\label{sec:uccd6_def}
% ═══════════════════════════════════════════════════════════════════════════════

中心差分系である基底 CCD（§\ref{sec:ccd_def}）を対流項に適用すると，
最高波数 $\xi=\pi$ における散逸ゼロから\textbf{Gibbs 振動}が許容される．
{\UCCD[6]} は，Chu--Fan の $D_1^{\CCD}$ の
\textbf{6 次分散精度を保ったまま}
高次ハイパー粘性 $\sigma_6|a|h^7(-D_2^{\CCD})^4$ を
\textbf{演算子内部に埋め込む}風上派生であり，
節点中心のままモーメント移流の $L^2$ 安定性を得る．
散逸を外殻 post-filter として加える DCCD（§\ref{sec:dccd_derivation}）と
ともに\textbf{節点中心 dissipation family} を構成し
（DCCD は post-filter，UCCD6 は演算子内部散逸という対比），
面中心化を行う FCCD（§\ref{sec:fccd_def}）は別 family に属する．
周期格子では Crank--Nicolson 時間積分との結合で無条件 $L^2$ 安定となる
（§\ref{sec:time_int}）．

% -------------------------------------------------------------------------------
\subsubsection{散逸動機の階層：1 次風上から $k=4$ ハイパー粘性へ}
\label{sec:uccd6_dissipation_hierarchy}
% -------------------------------------------------------------------------------

UCCD6 のハイパー粘性次数 $k=4$ は，より低次の風上散逸の自然な昇格として位置づけられる．
線形移流 $\partial_t u + a\partial_x u = 0$ に対する 1 次風上差分の修正方程式は
\begin{equation}
  \partial_t u + a\,\partial_x u
  \;=\; \tfrac{|a|\,h}{2}\,\partial_x^{2} u + \Ord{h^{2}},
  \label{eq:uccd6_upwind_modified}
\end{equation}
すなわち一般形 $-h^{2k-1}\partial_x^{2k}$ 型散逸の $k=1$ 構成
（$\partial_x^{2}$ の 2 階散逸；係数 $|a|h/2$）を帯びる．
この $k=1$ 散逸は中心 CCD 系の純虚数シンボル（§\ref{sec:dccd_embedding}）への
散逸補完として有効だが，低波数領域も一様に減衰させ
主項 $D_1^{\CCD}$ の $\Ord{h^6}$ 精度を律速する．
UCCD6 は同じ「演算子内部散逸」枠組のまま $k$ を昇格させた風上派生であり，
$k=4$（$h^{7}\partial_x^{8}$）が主項精度を保つ最小選択となる
（$k$ の必然性は次小節で精査）．
基底 CCD が既に解く $\omega_2$ シンボル（§\ref{sec:ccd_omega_weights}）を
$(-D_2^{\CCD})^{4}$ として再利用することで，
新たな高次差分テンプレートを追加せずに
Nyquist 選択散逸を演算子内部に埋め込む節点中心構成が得られる．

% -------------------------------------------------------------------------------
\subsubsection{演算子定義と符号}
\label{sec:uccd6_hyperviscosity}
% -------------------------------------------------------------------------------

線形移流 $\partial_t u + a\,\partial_x u = 0$ に対する半離散演算子：
\begin{equation}
  \partial_t U_j
  \;=\; L_h U_j
  \;=\; -a\,(D_1^{\CCD}U)_j \;-\; \sigma_6|a|h^{7}\,\bigl((-D_2^{\CCD})^{4}U\bigr)_j,
  \qquad \sigma_6 > 0.
  \label{eq:uccd6_semidiscrete}
\end{equation}
ここで $D_1^{\CCD}, D_2^{\CCD}$ は Chu--Fan CCD の 1 次・2 次導関数演算子
（§\ref{sec:ccd_def}）．第二項が\textbf{ハイパー粘性}であり，
$(-D_2^{\CCD})$ が正定値（§\ref{sec:ccd_omega_weights}）であることから
その 4 乗も正定値で，式~\eqref{eq:uccd6_semidiscrete} は\textbf{散逸演算子}を加えた形となる．

指数 7 の選択理由を確認する．周期格子での Fourier シンボル：
\begin{equation}
  \lambda(\theta)
  \;=\; -i\,\frac{a\,\omega_1(\theta)}{h}
  \;-\; \sigma_6|a|\,\frac{\omega_2(\theta)^{8}}{h},
  \qquad \theta = kh,
  \label{eq:uccd6_symbol}
\end{equation}
を用いると，$\omega_2^2 = \theta^2 + \Ord{\theta^8}$
（式~\eqref{eq:ccd_omega_12} 導出系）より
ハイパー粘性は物理波数で
$-\sigma_6|a|h^{7}\partial_x^{8}u$ として作用し，
主項 $D_1^{\CCD}$ の $\Ord{h^6}$ 切断誤差より\textbf{漸近的に劣位}．
指数 5（$h^{5}(-D_2^{\CCD})^3$）では散逸項が $\omega_2^6/h = \theta^6/h + \Ord{\theta^{12}/h}$
として作用し，主項 $D_1^{\CCD}$ の $\Ord{h^6}$ 切断誤差（$\theta - \theta^7/9450 + \Ord{\theta^9}$）と
低波数で同オーダー競合し精度を削る；
一方で指数 9（$h^{9}(-D_2^{\CCD})^5$）も低波数精度は保てるが，
同じ Nyquist 減衰オーダーを得るために余分な高階演算子を導入し，
境界閉包と丸め誤差への感度を増やす．
したがって指数 7・ハイパー粘性 4 乗は\textbf{次数保存の最小構成}である．

% -------------------------------------------------------------------------------
\subsubsection{半離散 $L^2$ 安定性（エネルギー法）}
\label{sec:uccd6_l2_stability}
% -------------------------------------------------------------------------------

周期格子の $\ell^2$ 内積 $\langle\cdot,\cdot\rangle_h$ の下で
$D_1^{\CCD}$ は歪 Hermite，$D_2^{\CCD}$ は Hermite かつ負定値．従って：
\begin{equation}
  \frac{1}{2}\frac{\mathrm{d}}{\mathrm{d}t}\|U\|_h^{2}
  \;=\; \langle U, \dot U\rangle_h
  \;=\; -\sigma_6|a|h^{7}\,\bigl\|(-D_2^{\CCD})^{2}U\bigr\|_h^{2}
  \;\le\; 0.
  \label{eq:uccd6_energy_id}
\end{equation}
移流項 $-a\langle U, D_1^{\CCD}U\rangle_h$ は歪 Hermite 性より消滅；
ハイパー粘性項が\textbf{厳密な散逸同定}を与える．
従って UCCD6 は全ての $\theta\in(0,\pi]$ で
モードごとに散逸する（$\omega_2(\theta)\ne 0$）．

% -------------------------------------------------------------------------------
\subsubsection{Crank--Nicolson による無条件安定（周期格子）}
\label{sec:uccd6_cn_stability}
% -------------------------------------------------------------------------------

時間刻み $\tau$ の Crank--Nicolson 離散化
$(I - \tfrac{\tau}{2}L_h)\,U^{n+1} = (I + \tfrac{\tau}{2}L_h)\,U^{n}$
の Fourier 増幅因子
$G(\theta) = (1+\tau\lambda(\theta)/2)/(1-\tau\lambda(\theta)/2)$
は，$\Re\lambda(\theta)\le 0$（式~\eqref{eq:uccd6_symbol}）より
$\tau, h$ によらず $|G(\theta)|\le 1$ を満たす；
すなわち UCCD6+CN は周期格子 $\ell^2$ 内積で\textbf{無条件安定}．
増幅因子の代数展開と一般 $L_h$ への拡張は §\ref{sec:time_int} で扱い，
有界域での GKS 解析は §\ref{sec:uccd6_boundary} 末で論じる．

\begin{remark}[陽解法の CFL 拘束]
\label{rem:uccd6_explicit_cfl}
TVD-RK3 などの陽解法を用いると $\omega_2^8$ 依存が $\theta=\pi$ 近傍で
$\tau_{\max}\le \Ord{h/(\sigma_6|a|\,\omega_{2,\max}^8)}$
（$\sigma_6$：UCCD6 散逸係数，$|a|$：移流速度の最大絶対値，
$\omega_{2,\max} := \max_{\theta\in[0,\pi]}\omega_2(\theta) = \omega_2(\pi)$ は
CCD 2 階微分の最大有効波数；式~\eqref{eq:ccd_omega_12} の $\theta=\pi$ 評価値）の
厳しい拘束を課す．中程度の $\sigma_6$ でも実質使えず，CN との組合せが必然．
\end{remark}

% -------------------------------------------------------------------------------
\subsubsection{境界閉包と GKS 文脈}
\label{sec:uccd6_boundary}
% -------------------------------------------------------------------------------

有界区間 $[0,L]$ では，周期格子で用いた
$D_1^{\CCD}$ の歪 Hermite 性と $D_2^{\CCD}$ の負定値性は
境界閉包（§\ref{sec:ccd_bc}）の選択に依存する．
したがって本節で証明する無条件安定性は周期格子の定理に限定し，
有界域での UCCD6 は境界閉包を含む行列スペクトルと MMS 収束で確認する
条件付きの設計要素として扱う．
$(-D_2^{\CCD})$ の 4 回適用は境界近傍で低次閉包の影響を受けるが，
$h^7$ 係数により低波数の主精度項とは分離される．
この境界影響は §\ref{sec:verify_ccd_convergence} の U1-d
（UCCD6 演算子検証）で測定し，
時間積分子としての CN 安定性は §\ref{sec:verify_time} の U8-c で検証する．
完全な Gustafsson--Kreiss--Sundstr\"om 解析は本稿の主張範囲に含めない．

% -------------------------------------------------------------------------------
\subsubsection{3 スキームの役割分担}
\label{sec:uccd6_scheme_comparison}
% -------------------------------------------------------------------------------

\begin{center}
\begin{tabular}{l l l l}
\hline
\textbf{スキーム} & \textbf{配置} & \textbf{散逸チャンネル} & \textbf{微分演算子精度}$^\star$ \\ \hline
Chu--Fan CCD（§\ref{sec:ccd_def}） & 節点 & なし & $\Ord{h^6}$ \\
DCCD（§\ref{sec:dccd_derivation}） & 節点 & 外殻 post-filter & $\Ord{h^6}$ \\
FCCD（§\ref{sec:fccd_def}） & \textbf{面} & なし（散逸ゼロ） & $\Ord{H^4}$ \\
\textbf{UCCD6（本節）} & 節点 & \textbf{演算子内部} & $\Ord{h^6}$ \\
\hline
\end{tabular}
\end{center}

{\footnotesize $^\star$ 微分演算子の formal 精度．DCCD では post-filter による移流フィルタ誤差 $\Ord{\varepsilon_d^\text{adv}\,h^2}$ が別途付随し，移流の実効精度は $\Ord{h^2}$ に律速される（§\ref{sec:dccd_filter_theory} 参照）．UCCD6 は $\omega_2^8$ 重みで Nyquist 帯域のみ選択減衰し bulk $\Ord{h^6}$ を維持．}

\noindent
\textbf{DCCD との数値特性比較．}
$\omega_2^8$ の選択的 Nyquist 散逸は DCCD の一様減衰 $H(\pi)=1-4\varepsilon_d$ より
dispersion preservation に優り，bulk 高周波の保持性で上位に位置する；
一方 DCCD の一様減衰は急峻 $C^0$ プロファイルのスペクトル振動を確実に抑制する．
このため両スキームは入力場の滑らかさに応じて使い分けられる
（具体的な適用場面は §~\ref{sec:ccd_summary} の比較表に集約）．

\noindent
\textbf{FCCD（§\ref{sec:fccd_def}）との関係．}
本章の派生スキームは出力ロカスで二つの family に分かれる：
節点中心 family（CCD / DCCD / UCCD6）と面中心 family（FCCD）である．
UCCD6 は節点中心 family の第 3 員で，演算子内部散逸により Gibbs を抑制する；
FCCD は面位置に値を持つ高次補間で $\Ord{H^4}$ を保ち，
\textbf{面ジェット}（§\ref{sec:face_jet_def}）を経由して節点中心 family の出力にも面 API を提供する．
両 family は同じ基底 CCD のシンボル $\omega_1, \omega_2$ から派生する．

% ═══════════════════════════════════════════════════════════════════════════════
